% =============================================================================
% Appendix A -- AI Usage Disclosure
% Required by Nova SBE Generative AI guidelines:
%   - AI used to RETRIEVE info: cite as personal communication, in-text
%   - AI used to GENERATE content: name the tool, explain method, date
%   - AI used as EDITING assistance: no disclosure needed
% Source: Nova SBE Guidelines Regarding Usage of Generative AI (current).
% =============================================================================

\section{AI Usage Disclosure}
\label{app:ai}

This appendix discloses the use of generative AI tools during the preparation
of this thesis, following the citation policy and the \emph{Know, Inform,
Assist} principles set out in the Nova SBE Generative AI guidelines
\citep{novasbe-genai}. A clarification is in order before the itemised list.
The thesis studies AI systems and produces AI artefacts: a frozen
Llama-3-8B-Instruct base \citep{llama3}, per-user LoRA adapters
\citep{hu2021lora} trained under DPO, IPO, SimPO and KTO, and the in-context
learning baselines. Those artefacts are the object of study and are documented
in \cref{sec:methodology,sec:design,sec:results}. The present appendix concerns
something different: the AI assistants used \emph{around} that work, as
research, coding and writing aids. Responsibility for every claim, figure and
line of code in the thesis remains solely with the author.

\subsection*{Tools used}

The assistants consulted during the project were Claude Sonnet 4.6, Claude code and Overleaf Writing Assistant.
Access took place between April 6 and May 20.
No paid third-party plug-ins or browsing connectors were used to send any
private or restricted data to these services; see the confidentiality note at
the end of this appendix.

\subsection*{AI used to retrieve information}

The Nova SBE guidelines treat AI as a source of \emph{personal communication}
when it is asked a question and returns an answer that is not independently
recoverable \citep{novasbe-genai}. Used in this way, AI assistants helped
clarify concepts whose primary sources are cited in the body of the thesis.
The main topics on which the assistants were queried are: the mechanics and
relative trade-offs of DPO, IPO, SimPO and KTO
\citep{rafailov2023dpo,azar2024ipo,meng2024simpo,ethayarajh2024kto}, the
geometry and hyper-parameter conventions of LoRA \citep{hu2021lora},
the interpretation of the paired Wilcoxon signed-rank test, and the operational
details of multi-adapter serving \citep{sheng2024slora}. In every case the
authoritative source consulted was the cited publication; the assistant's
output served as a comprehension aid and was discarded once the relevant
passage in the paper was located. Where a factual claim in the thesis rests on
one of these clarifications without an independent secondary source, that
claim was re-verified against the original publication before inclusion. No
in-text \emph{personal communication} citation appears in the body of the
thesis because no claim depends on AI output as a primary source.

\subsection*{AI used to generate content}

Per the Nova SBE guidelines, when AI is used to generate ideas or to develop a
plan, the tool, the method of use and the date of access must be acknowledged
\citep{novasbe-genai}. The following uses fall into this category.

\paragraph{Code and infrastructure.}
Portions of the training and evaluation pipeline were drafted with AI
assistance and then reviewed, executed and modified by the author. This
includes the per-user LoRA training loop, the YAML sweep configurations
referenced in \cref{sec:design}, the KTO row-expansion utility described in
\cref{sec:pair-gen}, the scoring code that groups KTO outputs back into pairs
by \texttt{pair\_id}, and the plotting code used to produce the figures in
\cref{sec:results}. Every code artefact was executed end-to-end on the actual
data.

\paragraph{LaTeX scaffolding.}
The section skeletons, the page-budget comment blocks at the top of each
section, and an initial draft of the \texttt{biblatex} entries in
\texttt{references.bib} were produced with AI assistance. Every bibliographic
entry was then checked against the published source; entries that could not be
verified are flagged in the \texttt{note} field of \texttt{references.bib} and
were either confirmed or removed before submission.

\paragraph{Ideation.} AI assistants were used as a sounding board for the framing of the cold-start problem in \cref{sec:intro}, the structure of the sweep grid in \cref{sec:exp}, and the business-case narrative in \cref{sec:business}. None of the analytical conclusions in those sections was adopted from an assistant's output verbatim: each was either re-derived from the experimental results or supported by a cited source.

\subsection*{AI used in editing assistance}

The Nova SBE guidelines exempt editing assistance from formal acknowledgement
\citep{novasbe-genai}. For transparency, the author notes that AI tools were
used for grammar, clarity and LaTeX-syntax suggestions on prose that the
author had already drafted.

\subsection*{Hallucination and verification}

Generative AI systems can produce confident but inaccurate output, including
fabricated citations \citep{novasbe-genai}. Three concrete safeguards were
applied throughout the project. First, every reference cited in the body of
the thesis was confirmed against the actual publication (DOI or arXiv ID);
items that could not be confirmed are tagged with a \texttt{Verify} note in
\texttt{references.bib} and were resolved before submission. Second, every
piece of AI-assisted code was executed on the real data and its output
compared against expected behaviour described in \cref{sec:methodology}, not
accepted on the strength of the assistant's explanation. Third, numerical
results, statistical tests and figure values reported in \cref{sec:results}
were re-derived from the raw sweep outputs by the author and are
reproducible from the artefacts kept alongside the thesis.

\subsection*{Confidentiality, copyright and data handling}

Both datasets used in this work are publicly released benchmarks: the Netflix
Prize dataset and the GoodReads-10k dataset, as described in
\cref{sec:data}. No personal data, no confidential material from any third
party, no NDA-covered content and no proprietary code base was uploaded to any
external AI service. Where snippets of the author's own code or LaTeX were
shared with AI assistants for editing assistance, the content involved no
material under embargo and no identifying information about the dataset
subjects.

\subsection*{Accountability}

In line with the Nova SBE guidelines, the responsibility to discuss, defend
and justify the merit and sources of the content of this thesis rests solely
with the author \citep{novasbe-genai}. Every methodological choice, every
analytical step and every conclusion reported in the body of the document can
be re-derived from the materials cited and reproduced from the artefacts kept
alongside the submission. The use of AI tools described above accelerated
drafting and verification; it did not substitute for the underlying decisions,
which remain the author's own.